\documentclass[conference]{IEEEtran}

% ─── Packages ─────────────────────────────────────────────────────────────────
\usepackage[utf8]{inputenc}
\usepackage{graphicx}
\usepackage{booktabs}
\usepackage{multirow}
\usepackage{amsmath}
\usepackage{amssymb}
\usepackage{xcolor}
\usepackage{microtype}
\usepackage{float}
\usepackage{cite}
\usepackage{hyperref}

% ─── Macros ───────────────────────────────────────────────────────────────────
\newcommand{\csgc}{\textsc{CSGC}}
\newcommand{\lean}{\textsc{CSGC}$_{\text{Lean}}$}
\newcommand{\tele}{\textsc{CSGC}$_{\text{Tele}}$}
\newcommand{\PR}{\textit{Pivot Recall}}

% ══════════════════════════════════════════════════════════════════════════════
\begin{document}

\title{CSGC: Compressing Conversation History Without Losing Corrections---
  Entity-Aware Supersession Detection and the Pivot Recall Metric}

\author{%
  \IEEEauthorblockN{Utkarsh Aggarwal}
  \IEEEauthorblockA{Independent Researcher\\
    utkarshaggarwal06@gmail.com}
}

\maketitle

% ─── Abstract ─────────────────────────────────────────────────────────────────
\begin{abstract}
Modern LLM-powered applications accumulate extensive conversation histories
that must be compressed before being re-injected into limited context windows.
Current approaches optimise for average semantic overlap with the original
text, yet ignore a critical temporal phenomenon: users and assistants
routinely \emph{correct} earlier facts during a conversation, and a
compressor that retains obsolete versions can degrade response accuracy.

We introduce \textbf{Context-Salience Greedy Compression (\csgc{})}, a
deterministic, three-stage pipeline---Redundancy Collapse, Salience Scoring,
and Greedy Budget Selection---that explicitly detects and deprioritises
superseded messages while promoting their corrections.
We also introduce \textbf{Pivot Recall}, a novel evaluation metric that tests
whether a compressed output contains the \emph{corrected} version of a fact
rather than the one it replaced.

Evaluated on 99 long-context English conversations from WildChat-1M at a
30\,\% byte budget, \lean{} achieves 91.2\,\% semantic retention
(matching LLMLingua-2's 89.1\,\% and running \textbf{4.2$\times$ faster})
while achieving a Pivot Recall of \textbf{58.3\,\%}---compared to 50.4\,\%
for Recency-based truncation and \textbf{5.5\,\%} for
LLMLingua-2. Our results suggest that perplexity-guided token distillation
does not inherently preserve factual corrections, and that lightweight entity-tracking recovers
the majority of this structural signal without any LLM calls.
\end{abstract}

\begin{IEEEkeywords}
context window compression, prompt engineering, information retrieval,
entity extraction, conversational AI, supersession
\end{IEEEkeywords}

% ══════════════════════════════════════════════════════════════════════════════
\section{Introduction}
\label{sec:intro}
% ─────────────────────────────────────────────────────────────────────────────

\IEEEPARstart{L}{arge} Language Models deployed in conversational settings
must re-read the entire interaction history on every forward pass. This
creates a well-known trade-off: longer context yields better recall but
incurs higher API costs, increased time-to-first-token (TTFT), and degraded
reasoning performance when irrelevant content dilutes the
prompt---the so-called \emph{``lost in the middle''} effect~\cite{liu2023lostmiddle}.

The standard response has been prompt compression: pruning the context to a
fixed budget before each new turn. Three families of methods exist.
\emph{Naive truncation} (Recency, Lead+Tail) is $\mathcal{O}(1)$ but
destroys either early or late context. \emph{Extractive heuristics} (TF-IDF,
TextRank) score messages by term statistics and retain the top-scoring subset.
\emph{Token-level distillation} (LLMLingua~\cite{jiang2023llmlingua},
LLMLingua-2~\cite{pan2024llmlingua2}) passes the full context through a
small language model to predict per-token keep/discard masks.

All three families share a fundamental assumption: \emph{the importance of a
message is static}. This assumption breaks in conversational settings. When a
user writes ``Let's use PostgreSQL'' and later writes ``Actually, switch to
MySQL---Postgres doesn't support our auth plugin,'' the first message is not
merely less important; it can lead to factual inconsistencies if retained without
the correction. We call such update events \textbf{supersession pivots}.

Our central contribution is to show that this phenomenon is both common and
consequential. Across 99 real conversations we detect 1,010 supersession
pivot pairs. Existing compressors fare poorly: LLMLingua-2 passes only
5.5\,\% of Pivot Recall checks because distillation models optimise for
fluency, not temporal fact ordering. Simple Recency truncation does better
(50.4\,\%) but at the cost of severe early-context loss (43.7\,\%
retention at 10\,\% budget).

We propose \csgc{}, a lightweight pipeline that: (1)~\textbf{detects}
supersession events using entity-aware sliding-window cosine similarity over
\texttt{e5-small-v2} embeddings~\cite{wang2024e5}; (2)~\textbf{scores}
retained messages via a weighted combination of TF-IDF, entity density, and
recency; and (3)~\textbf{selects} a byte-bounded subset via a greedy
knapsack, applying a $0.5\times$ penalty to superseded messages and a
$1.3\times$ boost to their corrections.

\csgc{} requires no GPU, makes no LLM API calls, and runs in under 800\,ms
per conversation---compatible with synchronous API middleware.

% ══════════════════════════════════════════════════════════════════════════════
\section{Related Work}
\label{sec:related}
% ─────────────────────────────────────────────────────────────────────────────

\subsection{Prompt Compression}
LLMLingua~\cite{jiang2023llmlingua} introduced selective context and
perplexity-guided token dropping. LLMLingua-2~\cite{pan2024llmlingua2} frames
compression as token-classification, training a 280M-parameter encoder on
GPT-4-distilled labels. RECOMP~\cite{xu2023recomp} and
AutoCompressor~\cite{chevalier2023autocompressor} extend these ideas to
retrieval-augmented generation. All of these methods operate at sub-word
granularity and cannot natively represent one message replacing another.

\subsection{Extractive Summarisation}
TextRank~\cite{mihalcea2004textrank} constructs a sentence-similarity graph
and applies PageRank to score sentences. TF-IDF variants treat each message
as a document and rank by corpus-relative term frequencies. Neither approach
models time nor the directed supersession relationship between messages.

\subsection{Context Management}
Liu et al.~\cite{liu2023lostmiddle} demonstrate empirically that LLMs attend
to the beginning and end of a long context far more than to the
middle---motivating the Lead+Tail baseline. Sliding window
attention~\cite{beltagy2020longformer} and memory
augmentation~\cite{wu2022memorizing} address this structurally but require
architectural changes. \csgc{} targets the orthogonal setting where the LLM
is a black-box API.

\subsection{Fact Supersession}
Current prompt compression literature primarily focuses on static context
summarisation and token distillation. The closest adjacent work studies fact
drift across training corpora rather than temporal updates within a single
conversation context~\cite{dhingra2022timedial}.

% ══════════════════════════════════════════════════════════════════════════════
\section{The \csgc{} Pipeline}
\label{sec:method}
% ─────────────────────────────────────────────────────────────────────────────

Let a conversation be a sequence of $n$ messages
$\mathcal{M} = \{m_1, \ldots, m_n\}$, each with role
$r_i \in \{\text{user},\text{assistant}\}$, textual content $t_i$, and byte
size $s_i$. Given a target byte budget $B$, \csgc{} produces a subsequence
$\mathcal{C} \subseteq \mathcal{M}$ such that $\sum_{m \in \mathcal{C}} s_m
\leq B$. The three stages are illustrated in Fig.~\ref{fig:arch}.

\begin{figure}[t]
  \centering
  % FIX: use your cleaned architecture image (jpg or pdf, both work)
  \includegraphics[width=\linewidth]{fig4_architecture.jpg}
  \caption{The \csgc{} three-stage pipeline. Stage~1 merges adjacent
    same-role messages and detects supersession pivots via entity-aware
    sliding-window similarity. Stage~2 assigns a composite salience score.
    Stage~3 greedily packs the highest-scoring messages into the byte budget.}
  \label{fig:arch}
\end{figure}

\subsection{Stage 1: Redundancy Collapse}
\label{ssec:collapse}

\textit{Same-role merging.}
Consecutive messages with identical roles are concatenated, collapsing reply
chains into a single logical unit.

\textit{Sliding-window deduplication.}
Let $\mathbf{e}_i \in \mathbb{R}^d$ be the $\ell_2$-normalised
\texttt{e5-small-v2} embedding of $t_i$. For each message $m_i$, we compute
cosine similarities against its $W\!=\!15$ predecessors:
\begin{equation}
  \text{sim}(i,j) = \mathbf{e}_i \cdot \mathbf{e}_j,\quad
  j \in \{\max(1,i{-}W),\ldots,i{-}1\}.
\end{equation}
If $\text{sim}(i,j) \geq \theta_\text{dup} = 0.95$, we compare entity sets
$E_i$ and $E_j$ extracted via spaCy with numeric/date regex fallbacks:
\begin{itemize}
  \item $E_i = E_j$ $\Rightarrow$ \textbf{Collapse}: $m_j$ is absorbed into
    $m_i$ (exact repetition).
  \item $E_i \neq E_j$ $\Rightarrow$ \textbf{Supersession}: $m_j$ is flagged
    \textsc{Obsolete}; $m_i$ is flagged \textsc{Correction}.
\end{itemize}

The key insight is that near-duplicate messages with matching entities
represent repetition, while diverging entities signal a factual update.

\subsection{Stage 2: Salience Scoring}
\label{ssec:salience}

Each message receives a composite salience score $\sigma_i \in [0,1]$:
\begin{equation}
  \sigma_i = w_\text{tfidf}\,\tilde{f}_i^{\text{tfidf}}
           + w_\text{ent}\,\tilde{f}_i^{\text{ent}}
           + w_\text{rec}\,\tilde{f}_i^{\text{rec}},
\end{equation}
where all components are min-max normalised within the conversation and
weights are $w_\text{tfidf}\!=\!0.50$, $w_\text{ent}\!=\!0.30$,
$w_\text{rec}\!=\!0.20$.

Supersession modifiers are applied post-scoring:
\begin{equation}
  \hat{\sigma}_i = \begin{cases}
    \sigma_i \times 0.5 & \text{if } m_i \in \textsc{Obsolete},\\
    \sigma_i \times 1.3 & \text{if } m_i \in \textsc{Correction},\\
    \sigma_i             & \text{otherwise.}
  \end{cases}
\end{equation}

These heuristic multipliers were selected to enforce strict precedence: a $0.5\times$ penalty aggressively down-weights obsolete content, while a $1.3\times$ boost elevates corrections above average salience to ensure they survive budget constraints.

\subsection{Stage 3: Greedy Budget Selection}
\label{ssec:selection}

Messages are sorted in descending $\hat{\sigma}_i$ order and greedily added
to $\mathcal{C}$ until budget $B$ is exhausted, then re-sorted
chronologically to preserve conversational coherence.

We evaluate two output variants:
\textbf{\lean{}} emits the full text of each selected message;
\textbf{\tele{}} additionally strips stop-words and non-entity filler tokens
via deterministic regex, achieving tighter compression at the cost of
grammaticality.

% ══════════════════════════════════════════════════════════════════════════════
\section{Pivot Recall: A New Evaluation Metric}
\label{sec:metric}
% ─────────────────────────────────────────────────────────────────────────────

Standard compression benchmarks measure average semantic similarity between
compressed and original context. A compressor can score 90\,\% semantic
retention while systematically discarding every factual correction, because
corrected and original facts share vocabulary and are locally similar.

\textit{Pivot QA set construction.}
We run Stage~1 over the initially sampled 100 conversations to extract supersession pairs
$\langle m_\text{old}, m_\text{new}\rangle$. (As noted, one conversation is subsequently filtered due to extraction failures). Across the 99 valid conversations, this yields \textbf{1,010 pivot QA pairs}. The gold target is the embedding
of $m_\text{new}$ (corrected fact); the distractor is the embedding of
$m_\text{old}$ (obsolete fact).

\textit{Scoring.}
For compressed output $\mathcal{C}$ and pivot pair $k$, let
$s_\text{new} = \max_{m \in \mathcal{C}}\,\mathbf{e}_m \cdot
\mathbf{e}_\text{new}^{(k)}$ and
$s_\text{old} = \max_{m \in \mathcal{C}}\,\mathbf{e}_m \cdot
\mathbf{e}_\text{old}^{(k)}$.
A pair passes iff:
\begin{equation}
  s_\text{new} \geq 0.92 \;\wedge\; s_\text{new} > s_\text{old}.
\end{equation}
\PR{} is the fraction of pairs that pass:
\begin{equation}
  \text{PR} = \frac{1}{K}\sum_{k=1}^{K}
  \mathbf{1}\!\left[s^{(k)}_\text{new} \geq 0.92
               \;\wedge\; s^{(k)}_\text{new} > s^{(k)}_\text{old}\right].
\end{equation}

The conjunctive condition is intentional: a method must retain the correction
\emph{more faithfully} than the error it supersedes.

% ══════════════════════════════════════════════════════════════════════════════
\section{Experimental Setup}
\label{sec:setup}
% ─────────────────────────────────────────────────────────────────────────────

\textit{Dataset.}
We sample 100 English conversations from WildChat-1M~\cite{zhao2024wildchat},
filtered for $\geq$10 turns. The random seed is fixed at 42. One conversation
was excluded due to an entity extraction failure, leaving 99 conversations for
the evaluation set (\csgc{} is parameter-free).

\textit{Baselines.}
\textbf{RAW} (uncompressed oracle);
\textbf{Lead+Tail} (first 15\,\%\,+\,last 15\,\%);
\textbf{Recency} (most recent messages up to budget);
\textbf{Random} (uniformly sampled messages);
\textbf{TF-IDF} (greedy selection by TF-IDF score);
\textbf{TextRank}~\cite{mihalcea2004textrank} (PageRank over similarity graph);
\textbf{LLMLingua-2}~\cite{pan2024llmlingua2} (run locally, CPU-only).

% FIX: Renamed from "General QA metric" to match actual implementation.
% The metric is purely e5-small-v2 cosine similarity; no LLM grading was used.
\textit{Semantic Retention metric.}
A gold assistant message from the first third of each conversation serves as a
semantic anchor. Using \texttt{e5-small-v2}~\cite{wang2024e5}, we compute the
maximum cosine similarity between the gold embedding and all messages in the
compressed output (pass threshold $\theta = 0.92$, score scaled $\times 100$
to report as a percentage).

\textit{Budget sweep.}
All methods are evaluated at 10, 20, 30, 40, 50, and 60\,\% byte budgets.
Results are averaged over the 99 valid conversations.

% ══════════════════════════════════════════════════════════════════════════════
\section{Results}
\label{sec:results}
% ─────────────────────────────────────────────────────────────────────────────

\subsection{Overall Benchmark at 30\,\% Budget}

Table~\ref{tab:main} and Fig.~\ref{fig:comp_vs_qa} report all methods at
the canonical 30\,\% budget.

\begin{table}[t]
  \centering
  \caption{Overall benchmark at 30\,\% byte budget (N\,=\,99 conversations).
    Methods sorted by semantic retention (Ret.\,\%). Runtime is
    per-conversation on a single CPU core with no GPU or LLM API calls.}
  \label{tab:main}
  \small
  \begin{tabular}{lrrrr}
    \toprule
    \textbf{Method} & \textbf{Comp.\,\%} & \textbf{Saved\,\%}
      & \textbf{Ret.\,\%} & \textbf{ms} \\
    \midrule
    RAW              & 100.0 &  0.0 & 99.9          &    0.0 \\
    Lead+Tail        &  27.5 & 72.5 & 96.6          &    0.0 \\
    TextRank         &  27.7 & 72.3 & 93.0          &   96.2 \\
    TF-IDF           &  28.8 & 71.2 & 93.0          &    9.8 \\
    Random           &  27.9 & 72.1 & 93.0          &    0.1 \\
    \lean{}          &  28.6 & 71.4 & \textbf{91.2} &  795.3 \\
    \tele{}          &  22.1 & 77.9 & 89.5          &  760.6 \\
    LLMLingua-2      &  34.4 & 65.6 & 89.1          & 3371.0 \\
    Recency          &  22.5 & 77.5 & 85.4          &    0.0 \\
    \bottomrule
  \end{tabular}
\end{table}

\begin{figure}[t]
  \centering
  \includegraphics[width=\linewidth]{fig1_SementicVsCompRatio.png}
  \caption{Compression ratio (lower = tighter) vs.\ semantic retention at
    30\,\% budget. \lean{} and LLMLingua-2 cluster near the efficient
    frontier; Recency is a visible outlier due to its recency bias.}
  \label{fig:comp_vs_qa}
\end{figure}

At 30\,\% budget, \lean{} achieves 91.2\,\% retention---2.1\,pp above
LLMLingua-2 (89.1\,\%) and 5.8\,pp above Recency (85.4\,\%). \tele{}
achieves stronger compression (22.1\,\% ratio vs.\ 28.6\,\% for \lean{})
while maintaining 89.5\,\% retention, outperforming LLMLingua-2 under a
stricter byte constraint.

The Random baseline (93.0\,\%) achieves competitive semantic retention because
WildChat conversations are topically coherent; any random sample retains
overlapping vocabulary with the gold reference. This underscores why Pivot
Recall is a necessary companion metric.

\subsection{Retention Across Budget Targets}

Table~\ref{tab:budget} and Fig.~\ref{fig:budget_curve} show semantic
retention across six budget levels.

\begin{table}[t]
  \centering
  \caption{Semantic retention (\%) across budget targets. Bold = best
    among compressed methods.}
  \label{tab:budget}
  \small
  \begin{tabular}{lrrrrrr}
    \toprule
    \textbf{Method} & \textbf{10} & \textbf{20} & \textbf{30}
      & \textbf{40} & \textbf{50} & \textbf{60} \\
    \midrule
    RAW         & 99.9 & 99.9 & 99.9 & 99.9 & 99.9 & 99.9 \\
    Lead+Tail   & \textbf{93.0} & \textbf{95.1} & \textbf{96.6}
                & \textbf{98.3} & \textbf{98.8} & \textbf{99.5} \\
    TF-IDF      & 89.7 & 92.0 & 93.0 & 95.3 & 95.4 & 96.2 \\
    TextRank    & 90.3 & 92.1 & 93.0 & 94.0 & 94.7 & 95.6 \\
    Random      & 88.9 & 91.1 & 93.0 & 93.9 & 95.0 & 96.2 \\
    \lean{}     & 88.5 & 90.7 & 91.2 & 93.3 & 94.0 & 95.0 \\
    \tele{}     & 87.2 & 89.2 & 89.5 & 90.9 & 91.4 & 92.2 \\
    LLMLingua-2 & 88.2 & 88.8 & 89.1 & 89.3 & 89.4 & 89.8 \\
    Recency     & 43.7 & 79.1 & 85.4 & 88.3 & 89.5 & 90.3 \\
    \bottomrule
  \end{tabular}
\end{table}

\begin{figure}[t]
  \centering
  \includegraphics[width=\linewidth]{fig2_QA_RetentionvsStorageBudget.png}
  \caption{Semantic retention vs.\ storage budget. Recency's sharp inflection
    between 10\,\%--20\,\% exposes its extreme sensitivity to context length.
    \lean{} scales consistently across the full range.}
  \label{fig:budget_curve}
\end{figure}

LLMLingua-2 exhibits a remarkably flat retention curve (88.2\,\% at 10\,\%,
89.8\,\% at 60\,\%), suggesting its token masks are largely budget-invariant.
\lean{} scales from 88.5\,\% to 95.0\,\%---a desirable property for systems
requiring fine-grained cost control. Recency's collapse at 10\,\% budget
(43.7\,\%) confirms the ``lost in the middle'' intuition.

\subsection{Pivot Recall: The Central Finding}

\begin{figure}[t]
  \centering
  \includegraphics[width=\linewidth]{Fig6_GenralQAvsPivotRecall.png}
  % FIX: caption corrected from "QA Accuracy" to "Semantic Retention"
  \caption{Semantic retention vs.\ Pivot Recall at 30\,\% budget (1,010
    supersession pairs). All methods achieve similar semantic retention
    (blue), yet LLMLingua-2 exhibits substantially lower Pivot Recall
    (5.5\,\%).}
  \label{fig:pivot}
\end{figure}

\begin{table}[t]
  \centering
  % FIX: renamed column header from "Gen. QA %" to "Sem. Ret. %"
  \caption{Semantic retention and Pivot Recall at 30\,\% budget. The
    \textit{Gap} column reveals how semantic retention alone overestimates
    fact-correction quality.}
  \label{tab:pivot}
  \small
  \begin{tabular}{lrrr}
    \toprule
    \textbf{Method} & \textbf{Sem. Ret.\,\%} & \textbf{Piv. Recall\,\%}
      & \textbf{Gap} \\
    \midrule
    RAW             & 99.9 & 91.3 &  8.6 \\
    TF-IDF          & 93.0 & 43.0 & 50.0 \\
    Recency         & 85.4 & 50.4 & 35.0 \\
    \lean{}         & 91.2 & 58.3 & \textbf{32.9} \\
    LLMLingua-2     & 89.1 &  5.5 & 83.6 \\
    \bottomrule
  \end{tabular}
\end{table}

\textit{LLMLingua-2 failure analysis.}
LLMLingua-2 achieves 89.1\,\% semantic retention but only 5.5\,\% Pivot
Recall---an 83.6\,pp gap. The cause is structural: LLMLingua-2's
token-classification objective rewards retaining fluent, high-perplexity
tokens. Factual corrections often consist of small numeric or proper-noun
substitutions (e.g., ``MySQL'' vs.\ ``PostgreSQL'') that are
\emph{low-perplexity in context} and thus preferentially dropped.
When original fact-bearing tokens survive the mask while correction tokens
are discarded, the compressed output provides incomplete context to the downstream LLM.

\textit{CSGC advantage.}
\lean{} achieves 58.3\,\% Pivot Recall, the smallest semantic-retention to
Pivot-Recall gap (32.9\,pp) among all compressed methods. The remaining gap
to the RAW oracle (91.3\,\%) reflects cases where budget pressure forces
omission of both old and new facts, or where the correcting message is too
short to outcompete longer, higher-salience messages.

\subsection{Computational Efficiency}

\begin{figure}[t]
  \centering
  \includegraphics[width=\linewidth]{fig3_runtime.png}
  \caption{Mean per-conversation runtime (log scale). \lean{} runs in
    795\,ms; LLMLingua-2 (3,371\,ms) is 4.2$\times$ slower. Both run
    entirely on CPU with no GPU or external API calls.}
  \label{fig:runtime}
\end{figure}

\lean{}'s 795\,ms runtime is dominated by the \texttt{e5-small-v2} embedding
pass in Stage~1. Replacing this with a BM25 or bag-of-words entity-overlap
approximation represents an immediate optimisation avenue. Even as
implemented, \lean{} comfortably fits the latency envelope of asynchronous
API middleware.

\subsection{Lean vs.\ Telegraphic}

\begin{figure}[t]
  \centering
  \includegraphics[width=\linewidth]{Fig5_leanvsTele.png}
  \caption{\lean{} vs.\ \tele{}: (left) semantic retention by budget target;
    (right) retention at matched actual compression ratios. \tele{}'s token
    pruning achieves an additional $\sim$6\,pp compression with a minor
    retention penalty.}
  \label{fig:leantele}
\end{figure}

\tele{} consistently achieves tighter compression (22.1\,\% vs.\ 28.6\,\%
ratio at 30\,\% budget) at a 1.7\,pp retention cost, making it the preferred
choice under extreme budget constraints.

% ══════════════════════════════════════════════════════════════════════════════
\section{Discussion}
\label{sec:discussion}
% ─────────────────────────────────────────────────────────────────────────────

\textit{Why Pivot Recall matters in practice.}
Consider an agentic coding assistant that, two hours into a session, corrects
its earlier recommendation from ``use asyncio'' to ``use threading---asyncio
doesn't support this library.'' A compressor that retains the original while
dropping the correction may cause the LLM to repeat outdated advice.
This failure mode is often not captured by average semantic retention metrics,
whereas Pivot Recall explicitly surfaces this class of errors.

\textit{On the competitiveness of simple baselines.}
Lead+Tail achieves 96.6\,\% semantic retention in under 0.1\,ms. For
applications where long-term coherence is unimportant and conversations are
short, Lead+Tail remains the rational choice. \csgc{} is most valuable in
long, multi-session conversations with multiple mid-course corrections---the
regime most common in agentic and enterprise chat deployments.

\textit{Limitations.}
(1)~\texttt{e5-small-v2} truncates inputs to 512 tokens; very long messages
may cause missed entity detections.
(2)~Entity extraction relies on spaCy NER and regex fallbacks; it cannot
detect corrections that do not surface recognisable entities (e.g., tonal or
reasoning-level shifts).
(3)~All numbers derive from WildChat-1M; generalisation to domain-specific
corpora (medical, legal, code) is untested.
(4)~The $0.5\times$ and $1.3\times$ modifiers are empirical heuristics designed
to prioritise corrections; future work should systematically ablate and optimise
these hyperparameters against Pivot Recall.

% ══════════════════════════════════════════════════════════════════════════════
\section{Conclusion}
\label{sec:conclusion}
% ─────────────────────────────────────────────────────────────────────────────

We presented \csgc{}, a deterministic prompt compression pipeline and the
first to explicitly model conversational supersession. Our central finding is
that relying solely on semantic retention may mask limitations in distillation-based
compressors: LLMLingua-2 achieves a Pivot Recall of only 5.5\,\% despite
89.1\,\% semantic retention, meaning it struggles to retain the updated factual
state of the conversation. \csgc{} narrows this gap to a
32.9\,pp spread by tracking entity-level supersession events and penalising
obsolete content at selection time, while running 4.2$\times$ faster than
LLMLingua-2 without any LLM calls.

We release the Pivot Recall metric and evaluation code to support future
benchmarking of temporally-aware compression methods, and encourage the
community to adopt Pivot Recall as a standard complement to semantic retention
in prompt compression evaluations.


% ─────────────────────────────────────────────────────────────────────────────
\begin{thebibliography}{99}

\bibitem{liu2023lostmiddle}
N.~F. Liu, K.~Lin, J.~Hewitt, A.~Paranjape, M.~Bevilacqua, F.~Petroni, and
  P.~Liang, ``Lost in the middle: How language models use long contexts,''
  \textit{arXiv preprint arXiv:2307.03172}, 2023.

\bibitem{jiang2023llmlingua}
H.~Jiang, Q.~Wu, C.-Y. Lin, Y.~Yang, and L.~Qiu, ``{LLMLingua}: Compressing
  prompts for accelerated inference of large language models,'' in
  \textit{Proc. Conf. Empirical Methods in Natural Language Processing
  (EMNLP)}, 2023, pp. 13358--13376.

\bibitem{pan2024llmlingua2}
Z.~Pan, Q.~Wu, H.~Jiang, M.~Xia, X.~Luo, J.~Zhang, Q.~Lin, V.~R\"{u}hle,
  Y.~Yang, C.-Y. Lin, H.~V. Zhao, L.~Qiu, and D.~Zhang, ``{LLMLingua-2}:
  Data distillation for efficient and faithful task-agnostic prompt
  compression,'' \textit{arXiv preprint arXiv:2403.12968}, 2024.

\bibitem{xu2023recomp}
F.~Xu, W.~Shi, and E.~Choi, ``{RECOMP}: Improving retrieval-augmented {LMs}
  with compression and selective augmentation,'' \textit{arXiv preprint
  arXiv:2310.04408}, 2023.

\bibitem{chevalier2023autocompressor}
A.~Chevalier, A.~Wettig, A.~Ajith, and D.~Chen, ``Adapting language models
  to compress contexts,'' \textit{arXiv preprint arXiv:2305.14788}, 2023.

\bibitem{mihalcea2004textrank}
R.~Mihalcea and P.~Tarau, ``{TextRank}: Bringing order into text,'' in
  \textit{Proc. Conf. Empirical Methods in Natural Language Processing
  (EMNLP)}, 2004, pp. 404--411.

\bibitem{beltagy2020longformer}
I.~Beltagy, M.~E. Peters, and A.~Cohan, ``Longformer: The long-document
  transformer,'' \textit{arXiv preprint arXiv:2004.05150}, 2020.

\bibitem{wu2022memorizing}
Y.~Wu, M.~N. Rabe, D.~Hutchins, and C.~Szegedy, ``Memorizing transformers,''
  in \textit{International Conference on Learning Representations (ICLR)},
  2022.

\bibitem{wang2024e5}
L.~Wang, N.~Yang, X.~Huang, L.~Yang, R.~Majumder, and F.~Wei, ``Multilingual
  {E5} text embeddings: A technical report,'' in \textit{Proc. 62nd Annual
  Meeting of the Association for Computational Linguistics (ACL)}, 2024,
  pp. 5356--5369.

\bibitem{zhao2024wildchat}
W.~Zhao, J.~T. Chiu, Y.~Deng, V.~I. Spitz, R.~Y. Pang, A.~M. Rush,
  Y.~Kim, and K.~Cho, ``{WildChat}: 1M {ChatGPT} interaction logs in the
  wild,'' \textit{arXiv preprint arXiv:2405.01470}, 2024.

\bibitem{dhingra2022timedial}
B.~Dhingra, J.~R. Cole, J.~M. Eisenschlos, D.~Gillick, J.~Eisenstein, and
  W.~W. Cohen, ``Time-aware language models as temporal knowledge bases,''
  \textit{Transactions of the Association for Computational Linguistics},
  vol.~10, pp. 257--273, 2022.

\end{thebibliography}
% ─────────────────────────────────────────────────────────────────────────────

\end{document}
